%%%%%%%%%%%%%%%%%%%%%%% file template.tex %%%%%%%%%%%%%%%%%%%%%%%%%
%
% This is a general template file for the LaTeX package SVJour3
% for Springer journals.          Springer Heidelberg 2010/09/16
%
% Copy it to a new file with a new name and use it as the basis
% for your article. Delete % signs as needed.
%
% This template includes a few options for different layouts and
% content for various journals. Please consult a previous issue of
% your journal as needed.
%
%%%%%%%%%%%%%%%%%%%%%%%%%%%%%%%%%%%%%%%%%%%%%%%%%%%%%%%%%%%%%%%%%%%
%
\RequirePackage{fix-cm}
%
%\documentclass{svjour3}                     % onecolumn (standard format)
\documentclass[smallcondensed]{svjour3}     % onecolumn (ditto)
%\documentclass[smallextended]{svjour3}       % onecolumn (second format)
%\documentclass[twocolumn]{svjour3}          % twocolumn
\journalname{Fire Technology}
%\documentclass[12pt]{article}
%\setlength\topmargin{0pt}
%\addtolength\topmargin{-\headheight}
%\addtolength\topmargin{-\headsep}
%\setlength\oddsidemargin{0pt}
%\setlength\textwidth{\paperwidth}
%\addtolength\textwidth{-2in}
%\setlength\textheight{\paperheight}
%\addtolength\textheight{-2in}
%\linespread{1.25}
\usepackage{layout}
%
\smartqed  % flush right qed marks, e.g. at end of proof
%
\usepackage{graphicx}
%\usepackage{subcaption}
%\usepackage{placeins}
\newcommand{\etal}{\textit{et al}. }

\graphicspath{{figs/}}

\begin{document}

\title{Wildland Fire Spread Modeling Using Convolutional Neural Networks
}
\author{Removed for Peer Review}

\institute{Removed for Peer Review \at }

\date{Received: date / Accepted: date}


\maketitle

\begin{abstract}
This paper presents a novel predictive analytics approach to estimating
the spread of a wildland fire using a convolutional neural network
(CNN). Simulated burn maps for use in this process were
generated at six hour intervals using the phenomological fire spread model
of Rothermel with 10,000 different combinations of input parameters.
The robustness of the approach is tested using 1,000 simulations not
included when training the CNN. Overall the predicted burn maps 
from the CNN-based approach agreed with simulation results,
with mean precision, sensitivity, and F-measure of 0.97, 0.92, and 0.93,
respectively. Noise in the input parameters was found to not significantly
impact the CNN-based predictions. The computational cost of the method
was found to be comparable to a phenomological model in homogenous spatial
conditions, and significantly better for heterogenous spatial conditions.
Although trained on predictions six hours apart, the
CNN-based approach is shown to be capable of predicting burn maps
further in the future by recursively using previous predictions as inputs to
the model. When the initial fire was small, the model tended to under-predict
fire spread; however, predictions generally improved as the fire grew.

\keywords{Wildland Fire \and  Machine Learning \and Neural Network \and Fire Spread \and Convolutional Neural Network \and CNN}
% \PACS{PACS code1 \and PACS code2 \and more}
% \subclass{MSC code1 \and MSC code2 \and more}
\end{abstract}




\section{Introduction}
\label{intro}

Wildland fire propagation is a complex process which involves the
interactions of many underlying physical phenomena. Since fully
resolving these processes remains a research effort researchers
have developed simplified models which describe the growth rate
of a wildland fire.
A new wildland fire spread model which is capable of predicting the 2-D
spatial-temporal flame front faster than real-time for diverse landscapes,
fuel types, and weather conditions is needed.
Existing wildland fire spread models generally fall into three
categories: stochastic, phenomological, and physical.
A number of different fire spread model approaches have been
proposed with each class of model having different advantages and limitations
\cite{weber1991modelling,sullivan2007a,sullivan2007b}.

Stochastic fire spread models are based on statistical analysis of
historical wildland fires and prescribed burns. Sullivan examined 14
stochastic models developed between 1990-2013 and found each could be
summarized by the functional form chosen to describe the impact of wind speed,
zero wind fire spread, and fuel moisture content \cite{sullivan2007b}.
The impact of wind on wildland fire spread generally follows a power law
where the exponent and pre-factor are related to the fuel type. The
impact of fuel moisture content on wildland fire spread is generally
modeled as strongly linear or weakly exponential based on the fuel type.
These models are favored by municipalities as they are capable of
providing rapid reasonable estimates of the overall fire risk of the
region for which they were derived. However, stochastic models fail when
predicting conditions outside the historical data, such as new weather
conditions, spatial regions, or different fuel types. In addition, these
models generally make no attempt at resolving the spatial-temporal flame
front of a wildland fire.

Phenomological fire spread models are based on conservation of
energy, but use experimental measurements to develop functional forms
rather than modeling from first principles \cite{weber1991modelling,sullivan2007b}.
The most widely
known phenomological fire spread model is the model of Rothermel
\cite{rothermel1972mathematical,scott2005standard} which uses
empirical correlations for heat source and sink terms \cite{weber1991modelling}.
The baseline rate of spread (with no wind or slope) is based on fuel density,
type, and moisture content. The impact of wind and slope is modeled as a multiplier
of the baseline rate of spread. Although developed to determine the rate of spread
from a single ignition source in a single direction, these models have been expanded
to work in 2-D using Huygens' principle \cite{finney1999mechanistic,finney1998farsite}
and the level set method \cite{rehm2009fire,lautenberger2013wildland}.
These models are favored by foresters and
firefighters in the field as they are capable of providing rapid estimates
of the rate of spread \cite{simeoni2015wildland}. However, phenomological models
inherit similar issues to stochastic models when predicting conditions
outside the historical data used in deriving the empirical models. In addition,
expanding the model to 2-D to predict the spatial-temporal flame front significantly
increases the computational requirements.
Researchers have used
reduced order modeling to decrease the computational cost of evaluating the
2-D rate of spread using phenomological models \cite{lattimer2016computationally}.
However, it is difficult to apply this framework to other models as new derivations
are required.

Physical fire spread models are based on the fundamental chemistry and physics of
combustion, and fire spread \cite{sullivan2007a}. Physical models can be further 
subdivided into simple and detailed models. Simple physical models such as that of
Weber consider only the transport of energy, neglecting calculation of the full
flow field for computational reasons \cite{weber1991modelling}. These models
are more general than stochastic or phenomological models and can be close to
real-time \cite{simeoni2015wildland}. Detailed physical models such as
Wildland-Urban Interface Fire Dynamics Simulator (WFDS) consider the transport of
mass, momentum, and energy in addition to combustion in a multiphase approach
through computational fluid dynamics (CFD) \cite{mell2007physics}. These models
are favored by researchers interested in the fundamental mechanisms of wildland
fire spread as the basis on first principles reduces the impact of user intuition
on the results. However, the computational cost to model moderate domains
($2.25km^{2}$) is significant \cite{mell2007physics}.

Predictive analytics uses simplified functions to represent underlying connections
between data to make estimates of future events, often using machine learning.
Several researchers have presented predictive analytics frameworks to estimate the
total burned area of a fire based on meteorological data 
\cite{safi2013prediction,castelli2015predicting,storer2016pso,naganathan2016wildfire,cao2017wildfire}.
Each of these methods is capable of predicting the total burned area; however none
attempts to estimate the spatial-temporal fire distribution of a wildland fire front.
A predictive analytics approach to predicting the spatially-resolved final burned
region of a wildland fire was presented by McCormick.
The model used an artificial neural network to classify
the the center pixel of a 3x3 neighborhood of pixels as burned or unburned
\cite{mccormick2001toward,mccormick2002developing}.
The initial results are promising; however, the model relies on the assumed
ellipsoidal growth profile in the direction of wind
\cite{finney1999mechanistic} and only predicts the final burned region but 
not the evolution of spread.

A method to predict fully resolved spatial-temporal burn maps is desired.
The fundamental principle which makes convolutional neural networks (CNNs)
versatile is the capability to learn how to represent complex
shapes as combinations of high level feature maps. Krizhevsky showed
many of the features learned by the CNN in the ImageNet competition
described the inter-relationship of the 3 color channels \cite{krizhevsky2012imagenet}.
As an analogy to image classification, data such as elevation, moisture content, and 
wind speed can be treated as channels in an image. Given enough data, a CNN will be
able to learn relationships between these physical parameters which can then be used
to predict a future burn map. 

The objective of this study is to apply a trained CNN framework to predict the
spatial-temporal distribution of a wildland fire front in homogenous vegetation
without relying on any other models at runtime. Data for use in
training and testing the network was generated using Rothermel's phenomological
model. The quantitative performance of the model on 1,000 simulations not included
during training is presented, and
the sensitivity of the network to noise is examined.
The work presented herein demonstrates the concept
on a simple configuration with future work to expand the method to use
experimental data with heterogeneous spatial conditions.

\section{Methods}
\label{s:Methods}

A schematic showing a high level view of the solution algorithm is shown
in Fig.~\ref{fig:exampleIO}.
Each primary driver of wildland fire spread is included as a channel in an
image which is input to the CNN.
The CNN then uses its prior training to predict a new image with two channels
corresponding to the probability a pixel has burned or not burned in the future.
The two probability masks are then post-processed to output a single future
burn map. The temporal resolution of six hours and spatial resolution of 1 pixel/km
physically correspond to available global satellite measurements.
The following subsections describe the simulation conditions, network architecture,
post-processing, and performance metrics used in this work.

% For two-column wide figures use
\begin{figure*}[htb]
% Use the relevant command to insert your figure file.
% For example, with the graphicx package use
  \includegraphics[width=0.99\textwidth]{Fig1.eps}
% figure caption is below the figure
\caption{Schematic of solution algorithm. The left set of images show the
different channels used as inputs to the neural network. The values for each
data channel are colorized based on the values shown in
Table~\ref{tab:paramsLimits} and Table~\ref{tab:paramsExample}.}
\label{fig:exampleIO}       % Give a unique label
\end{figure*}



\subsection{Wildland Fire Prediction}
\label{ss:Wfp}

Data for this study was generated using the surface fire spread model
presented by Rothermel/Albini
\cite{rothermel1972mathematical,scott2005standard,albini1976estimating}.
In the Rothermel/Alibini phenomological fire spread model, the peak
surface fire spread rate, $V_{s,peak}$ is calculated using the equation
\begin{equation}
V_{s,peak} = \frac{Q''\zeta}{\rho\epsilon Q_{ig}}\left(1+\phi_{s}+\phi_{w}\right)
\label{eq:rothermel}
\end{equation}
where $Q''$ is the heat release rate per unit area, $\rho$ is the
fuel density, $Q_{ig}$ is the heat of pre-ignition, $\zeta$ is the
propagating flux ratio (percentage of heat released which pre-ignites
fuel), $\epsilon$ is the effective heating number (percentage of fuel
which is involved in ignition), $\phi_{w}$ is the wind coefficient,
and $\phi_{s}$ is the slope coefficient.

Various researchers have developed empirical relationships for the
different parameters in Eq.~\ref{eq:rothermel}. A commonly used approach
in the literature is to specify $Q''$, $\rho$, $\zeta$, $\epsilon$
based on classifying the primary fuel in a region into a fuel model.
A total of 53 fuel models were considered in this work including
13 developed by Rothermal/Albini \cite{rothermel1972mathematical,albini1976estimating},
and 40 developed by Scott \cite{scott2005standard}. Rothermel presented
an empirical relationship for $Q_{ig}$ based on the fuel model and
moisture content, and Scott extended the relationship to handle dynamic
fuel models. Rothermel presented empirical relationships for $\phi_{w}$
and $\phi_{s}$ based on fuel model, midflame wind speed, and slope.
Andrews presented an algorithm to adjust typical atmospheric wind
measurements (10m or 20ft) to mindflame wind speed based on three
additional parameters describing the upper story vegetation
(canopy cover, canopy height, and crown ratio)
\cite{andrews2012modeling}. Researchers have shown wildland fires grow
in a generally ellipsoidal shape for homogenous spatial conditions
based on $V_{s,peak}$ and wind speed
\cite{finney1999mechanistic,wagner1969simple,green1983fire}.
The primary drivers in the fire spread model were identified as
landscape (slope, aspect, and fuel model type),
moisture content (1-hour, 10-hour, 100-hour, live woody, and live herbaceous),
canopy type (height, ratio, and percent coverage),
and 10m wind (intensity and direction). The allowable bounds used for
each parameter in this work are shown in Table~\ref{tab:paramsLimits}.
The fuel model types were assigned indexes based on the peak rate of spread
under the same spread conditions (low moisture, 10 mph wind up a 0.5 slope).
Since slope and aspect can be summarized as a 2-D difference in elevation,
a single channel for elevation was used instead of two channels for slope
and aspect in the neural network.

% For tables use
\begin{table}[htb]
\centering
% table caption is above the table
\caption{Limits of each parameter in study.}
\label{tab:paramsLimits}       % Give a unique label
% For LaTeX tables use
\begin{tabular}{llll}
\hline\noalign{\smallskip}
Parameter & Unit & Min & Max  \\
\noalign{\smallskip}\hline\noalign{\smallskip}
Aspect & Degrees & 0.0 & 360 \\
Fuel Model & Index & 0.0 & 53 \\
Slope & Fraction & 0.0 & 1.0 \\
\hline
1-Hr Moisture & Percent & 1.0 & 40 \\
10-Hr Moisture & Percent & 1.0 & 40 \\
100-Hr Moisture & Percent & 1.0 & 40 \\
Live Herbaceous Moisture & Percent & 30 & 100 \\
Live Woody Moisture & Percent & 30 & 100 \\
\hline
Canopy Cover & Percent & 0.0 & 1.0 \\
Canopy Height & Feet & 1.0 & 20 \\
Crown Ratio & Fraction & 0.1 & 1.0 \\
\hline
Wind Direction & Degrees & 0.0 & 360 \\
Wind Velocity & Mi/Hr & 0.0 & 30 \\
\noalign{\smallskip}\hline
\end{tabular}
\end{table}

A custom implementation of Rothermel's fire spread model was developed
in Python to streamline the integration with the CNN software.
The implemention of Rothermel's model in BehavePlus was used to validate
the simulation framework \cite{andrews2009behaveplus}. A total of
10,000 different combinations of the parameters shown in Table~\ref{tab:paramsLimits}
were simulated. The total time to run 10,000 simulations was 30.6 seconds.
Raster images of the burned map were generated at a resolution of
1 pixel/km every six hours for a total of 24 hours.
Three cases for use in the CNN were generated from each simulation
by creating pairs of burned maps six hours apart. The zero hour to six hour
pair was not considered as a case for the CNN as the zero hour burned image
always contains no fire.
Example burned maps from one simulation are shown in
Fig.~\ref{fig:firePerimeters} for the parameter values
shown in Table~\ref{tab:paramsExample}.

\begin{figure}[htb]
\centering
% Use the relevant command to insert your figure file.
% For example, with the graphicx package use
  \includegraphics[width=0.45\textwidth]{Fig2.eps}
% figure caption is below the figure
\caption{Example Simulated Burn Maps}
\label{fig:firePerimeters}       % Give a unique label
\end{figure}

\begin{table}[htb]
\centering
% table caption is above the table
\caption{Parameters for Example Simulation}
\label{tab:paramsExample}       % Give a unique label
% For LaTeX tables use
\begin{tabular}{llll}
\hline\noalign{\smallskip}
Parameter & Unit & Value \\
\noalign{\smallskip}\hline\noalign{\smallskip}
Aspect & Degrees & 130 \\
Fuel Model & Index & FM1 (44) \\
Slope & Fraction & 0.8 \\
\hline
1-Hr Moisture & Percent & 5.3 \\
10-Hr Moisture & Percent & 6.3 \\
100-Hr Moisture & Percent & 7.3 \\
Live Herbaceous Moisture & Percent & 69 \\
Live Woody Moisture & Percent & 49 \\
\hline
Canopy Cover & Percent & 0.7 \\
Canopy Height & Feet & 14 \\
Crown Ratio & Fraction & 0.2 \\
\hline
Wind Direction & Degrees & 34 \\
Wind Velocity & Miles per Hour & 13.5 \\
\noalign{\smallskip}\hline
\end{tabular}
\end{table}
















\subsection{Network Architecture}
\label{ss:Na}


At a fundamental level, artificial neural networks are massively parallel
equations which have the capability to store observed knowledge about a
problem to make predictions of new inputs. A convolutional neural network
(CNN) assumes the input data has distinct
spatial dependence within the input parameters. Since the network
assumes spatial dependence, less connections need to be made to
inputs which are far from each other. This allows a CNN to contain much
fewer connections and parameters than a similarly sized standard feed
forward network with minimal loss in optimal performance for appropriate
problems. In addition, this makes it possible to use deeper and more broad hidden
layers without increasing computational requirements beyond what is
feasible on current technology \cite{krizhevsky2012imagenet}.
Representing the spread of a wildland fire front with a CNN is
reasonable as wildland fire spread is a local phenomena \cite{finney1999mechanistic}.

The CNN architecture used in this work is shown in Fig.~\ref{fig:cnnArchitecture}.
The input images were 50x50 pixels with 13 image channels corresponding to the
image stack shown in Fig.~\ref{fig:exampleIO}. The output image contains 50x50 pixels
with two image channels corresponding to the probability the burn map has reached
a pixel and the probability the burn map has not reached a pixel.
A total of seven hidden layers
are included in the network including two convolutional, two max pooling, one fully connected
classification, one up-convolutional, and one soft max layers. The number of filters
and step size in each convolutional and up-convolutional layers and number of neurons
in fully connected layers were specified to steadily decrease the degrees of freedom from the
32,500 in the input layer (50x50x13 inputs) to the desired degrees of freedom of
5,000 in the output layer (50x50x2 probability of fire and not fire).

% For two-column wide figures use
\begin{figure*}[htb!]
\centering
% Use the relevant command to insert your figure file.
% For example, with the graphicx package use
  \includegraphics[width=0.95\textwidth]{Fig3.eps}
% figure caption is below the figure
\caption{Convolutional Neural Network Architecture.}
\label{fig:cnnArchitecture}       % Give a unique label
\end{figure*}

All layers in the network presented in this work used leaky rectified linear
unit (ReLU) activation functions except the fully connected layer which used
hyperbolic tangent (TanH) activation functions and the output layer which used
a softmax activation function.
Researchers have shown deep neural networks learn more quickly and more
accurately using leaky or parametric ReLU than logistic or TanH for activation
\cite{maas2013rectifier,he2015delving}.
However, since the activation is unbounded, ReLU is not appropriate for
classification layers.
Since the intent of the fully connected layer is to classify the feature maps into
high level descriptions of the fire and environment, a TanH activation function
is appropriate.
Since the desired output is a probabilistic estimate of whether or not each
pixel would be burned, a soft max activation function is appropriate to scale the
activations of the fire and non-fire probability masks.


The network architecture was built using the Python 3 bindings for TensorFlow
\cite{tensorflow2015-whitepaper}. The models were trained using stochastic
gradient descent with a batch size of 100 samples. All weights and biases were
initialized from a uniform distribution between -1 and 1. The learning rate was
fixed for all layers throughout training at 0.0001.
The cost function used in training was based on sum square error.
Over-fitting was reduced by using 50\% dropout on the input layer and shuffling
the order of the samples during training.
The network was trained using 9,000 simulations (27,000 pairs of burn maps)
for 50,000 cycles
using a single NVIDIA Quadro K620. The total time to train the network was
18 hours 7 minutes. The total time to evaluate the network with all 10,000
simulations (30,000 pairs of burn maps) was 38.1 seconds.





\subsection{Performance Metrics}
\label{ss:PM}

The metrics used to quantify the performance of the CNN in this study
were precision, sensitivity, and F-measure. For each metric, the range
of possible values is zero to one, with a perfect score being one.
The precision, $P$, is a measure of commission errors (predicting a
fire where there was not fire) and is defined as
\begin{equation}
P = \frac{t_{p}}{t_{p}+f_{p}}
\label{eq:precision}
\end{equation}
where $t_{p}$ is the number of correctly identified fire pixels, and
$f_{p}$ is the number of falsely identified fire pixels.
The sensitivity, $S$, is a measure of omission errors (predicting no fire where there was a fire) and is defined as
\begin{equation}
S = \frac{t_{p}}{t_{p}+f_{n}}
\label{eq:sensitivity}
\end{equation}
where $f_{n}$ is the number of fire pixels which were identified
as non-fire pixels.
F-measure, $F$, is an overall measure of performance defined
as the harmonic mean of $P$ and $S$,
\begin{equation}
F = 2\cdot\frac{P\cdot S}{P+S}.
\label{eq:fmeasure}
\end{equation}



%\FloatBarrier

\subsection{Post Processing}
\label{ss:PP}

The output layer of the CNN contains two normalized probability masks, one for
fire and one for not fire. The normalized probability mask for fire is
post-processed to convert the probabilistic estimate of burned or unburned to
a single contour. A 3x3 median filter is applied to smooth the image. A threshold value
on probability of fire is used to determine whether or not each pixel has been
burned. 

The optimal threshold to use in post-processing was determined by calculating
the mean F-measure for CNN predictions of the 9,000 training simulations (27,000
pairs of burn maps) with thresholds ranging from 0.01 to 0.99.
The mean F-measure was found to be mostly independent of the post-processing threshold
in the range of 0.2-0.6 as shown in Fig.~\ref{fig:fMeasureVsThreshold}.
The maximum mean F-measure of the training data was calculated with a threshold
of 0.41. This threshold was fixed and used for post-processing throughout
this work. An example neural network prediction before and after post-processing
is shown in Fig.~\ref{fig:postProcess}.

\begin{figure}[htbp]
\centering
% Use the relevant command to insert your figure file.
% For example, with the graphicx package use
  \includegraphics[width=0.45\textwidth]{Fig4.eps}
% figure caption is below the figure
\caption{Mean F-Measure of CNN predictions of 9,000 training simulations (27,000 pairs of burn maps) for different threshold values.}
\label{fig:fMeasureVsThreshold}       % Give a unique label
\end{figure}

\begin{figure*}[htb]
	\centering
	\includegraphics[width=0.46\textwidth]{Fig5a.eps}
	~
	\includegraphics[width=0.46\textwidth]{Fig5b.eps}
	\\
	(a)\hspace{0.46\textwidth}(b)
\caption{Example CNN prediction of burn map probability (a) Directly predicted from CNN (b) After post-processing}
\label{fig:postProcess}       % Give a unique label
\end{figure*}











\section{Results}
\label{s:Results}

The robustness of the neural network to predict new fires was examined
by considering 1,000 simulations (3,000 pairs of burn maps) which were not
included when training the network.
Sample CNN predictions from five of these test cases are
compared with simulation predictions in Fig.~\ref{fig:exampleNetwork}.
Figure~\ref{fig:exampleNetwork}a shows the initial and next burn map
from the simulation.
Figure~\ref{fig:exampleNetwork}b shows the final burn map predicted by the CNN.
Figure~\ref{fig:exampleNetwork}c highlights pixels which the CNN prediction did not match
the simulation predictions. Pixels shown as black represent commission
errors (false positive of fire), and pixels shown as orange represent omission errors
(false negative of fire).

\begin{figure*}[htbp]
  \centering
	\includegraphics[height=0.17\textheight]{Fig6a1.eps}
	~
	\includegraphics[height=0.17\textheight]{Fig6b1.eps}
	~
	\includegraphics[height=0.17\textheight]{Fig6c1.eps}
	\\
	\includegraphics[height=0.17\textheight]{Fig6a2.eps}
	~
	\includegraphics[height=0.17\textheight]{Fig6b2.eps}
	~
	\includegraphics[height=0.17\textheight]{Fig6c2.eps}
	\\
	\includegraphics[height=0.17\textheight]{Fig6a3.eps}
	~
	\includegraphics[height=0.17\textheight]{Fig6b3.eps}
	~
	\includegraphics[height=0.17\textheight]{Fig6c3.eps}
	\\
	\includegraphics[height=0.17\textheight]{Fig6a4.eps}
	~
	\includegraphics[height=0.17\textheight]{Fig6b4.eps}
	~
	\includegraphics[height=0.17\textheight]{Fig6c4.eps}
	\\
	\includegraphics[height=0.17\textheight]{Fig6a5.eps}
	~
	\includegraphics[height=0.17\textheight]{Fig6b5.eps}
	~
	\includegraphics[height=0.17\textheight]{Fig6c5.eps}
	\\
\hspace{0.033\textwidth}(a)\hspace{0.29\textwidth}(b)\hspace{0.34\textwidth}(c)
\caption{Example CNN prediction results (a) Simulation burn maps (b) CNN predicted burn map (c) Classification error}
\label{fig:exampleNetwork}       % Give a unique label
\end{figure*}

The mean precision, sensitivity, and F-measure of the 1,000 test
simulations (3,000 burn map pairs) are shown in Table~\ref{tab:metrics}.
Each metric ranges from zero to one with a perfect score being one. 
The distribution of F-measure for the 3,000 burn map pairs are shown
in Fig.~\ref{fig:fMeasureDistribution}.
Percentiles of each metric are shown in Table~\ref{tab:metrics} to
quantify the spread of each score. Here a percentile is defined as
\begin{equation}
X = \frac{1}{N_{total}}\int_{Y_{min}}^{1}N(Y)dy
\label{eq:percentile}
\end{equation}
where $X$ is the percentile, $N(Y)$ is the number of burn maps achieving
a specific score, $N_{total}$ is the total burn maps $Y_{min}$ is the minimum
score. For example, the F-measure of 0.86 for $X=80\%$ shown in
Table~\ref{tab:metrics} means 80\% (2400/3000) of the burn map pairs
had an F-measure of 0.86 or higher.

\begin{table}[htb]
\centering
% table caption is above the table
\caption{Performance Metrics of CNN Predictions of Test Cases}
\label{tab:metrics}       % Give a unique label
% For LaTeX tables use
\begin{tabular*}{0.75\textwidth}{l @{\extracolsep{\fill}} cccc}
\hline\noalign{\smallskip}
Parameter & Mean & $X=80\%$ & $X=90\%$ & $X=95\%$\\
\noalign{\smallskip}\hline\noalign{\smallskip}
Precision & 0.97 & 0.95 & 0.85 & 0.79\\
Sensitivity & 0.92 & 0.80 & 0.67 & 0.59\\
F-Measure & 0.93 & 0.86 & 0.80 & 0.73\\
\noalign{\smallskip}\hline
\end{tabular*}
\end{table}

\begin{figure}[htbp]
\centering
% Use the relevant command to insert your figure file.
% For example, with the graphicx package use
  \includegraphics[width=0.45\textwidth]{Fig7.eps}
% figure caption is below the figure
\caption{F-Measure distribution of CNN predictions of test cases.}
\label{fig:fMeasureDistribution}       % Give a unique label
\end{figure}



















\section{Discussion}
\label{s:Discussion}

The overall shape of the burn maps predicted by the CNN are consistent
with the simulations for the 1,000 test simulations (3,000 burn map pairs)
examined in this work, as shown in the examples in Fig.~\ref{fig:exampleNetwork}.
The burn maps predicted by the CNN do not contain non-physical
holes or excessive noise. The direction of maximum growth is captured well.
Figure~\ref{fig:exampleNetwork} shows the CNN is able to predict the
growth of small ($1km^{2}$, rows 2 and 4), intermediate ($10km^{2}$, rows 1 and 5), and
large ($100km^{2}$ row 3) fires.

Across all test cases the was not significant bias
observed in over-predicting or under-predicting fires, as is shown by the
comparable mean sensitivity and precision shown in Table~\ref{tab:metrics}.
However, the spread in sensitivity is higher than the spread in precision
as is shown by $X=80\%$, $X=90\%$, and $X=95\%$ shown in Table~\ref{tab:metrics}.
Examining the histogram of $F$ shown in Fig.~\ref{fig:fMeasureDistribution},
there is a sharp drop in number of occurrences at $F < 0.8$.
The initial fire size of 82\% of the cases where $F < 0.8$ was found to be
nine pixels or less, as shown in Fig.~\ref{fig:fireSize}.
Combined with the percentiles shown in Table~\ref{tab:metrics}, this shows
the sensitivity is impacted to a greater extent than precision for small fire
sizes. Since a CNN relies on feature recognition, low feature density in
the inputs (such as fires less than 9 pixels in size) can lead to a decrease
in the accuracy of the model. If smaller fires are of primary interest,
this effect could be reduced by increasing the spatial resolution of the model.

\begin{figure}[htbp]
\centering
% Use the relevant command to insert your figure file.
% For example, with the graphicx package use
  \includegraphics[width=0.45\textwidth]{Fig8.eps}
% figure caption is below the figure
\caption{Distribution of input fire size for all cases where $F < 0.8$.}
\label{fig:fireSize}       % Give a unique label
\end{figure}

The sensitivity of the model to noise in the input parameters was examined
by adding noise to the 5 examples shown in Fig.~\ref{fig:exampleNetwork}
and examining the impact on $S$, $P$, and $F$. Noise was added by multiplying
each input channel by a random value from a log-normal distribution with a
mean of 0 and a standard deviation of 1. Note the input and output burn maps
were not changed. A total of 3,000 iterations for each of the five test
cases were predicted by the CNN. The zero noise, noise mean, and noise 80\%
scores for each case for each metric are shown in Table~\ref{tab:noise}.
The largest impact was observed in the last test case where the mean and 80\%
sensitivity dropped from the baseline score by 0.07 and 0.14,
respectively. This shows the addition of noise resulted in an increase in
the under prediction of the burn map.
However, the minimal overall impact on the mean and 80\% scores
shows the general shape of the predicted burn maps are still in agreement
with simulation results.

\begin{table}[htb]
\centering
% table caption is above the table
\caption{Impact of Noise on Performance Metrics of CNN Predictions of Test Cases}
\label{tab:noise}       % Give a unique label
% For LaTeX tables use
\begin{tabular*}{0.95\textwidth}{c @{\extracolsep{\fill}} cccccccc}
\multicolumn{3}{c}{Sensitivity} & \multicolumn{3}{c}{Precision} & \multicolumn{3}{c}{F-Measure} \\

Base & Noise & $80\%$ & Base & Noise & $80\%$ & Base & Noise & $80\%$ \\
\noalign{\smallskip}\hline\noalign{\smallskip}
0.99 & 0.98 & 0.99 & 0.98 & 0.98 & 0.98 & 0.98 & 0.98 & 0.98 \\
0.93 & 0.91 & 0.93 & 1.00 & 1.00 & 1.00 & 0.96 & 0.94 & 0.96 \\
0.98 & 0.97 & 0.95 & 0.99 & 0.99 & 0.99 & 0.98 & 0.98 & 0.97 \\
1.00 & 1.00 & 1.00 & 1.00 & 0.93 & 1.00 & 1.00 & 0.95 & 1.00 \\
0.98 & 0.91 & 0.84 & 0.92 & 0.90 & 0.85 & 0.95 & 0.90 & 0.88 \\
\noalign{\smallskip}\hline
\end{tabular*}
\end{table}

One of the key benefits of the CNN architecture presented in
this work is the scalability in heterogeneous spatial conditions.
Recall the total time to run 10,000 simulations using the model
of Rothermel was 30.6s compared to the 38.1s required for the CNN.
The phenomological model will likely perform faster whenever
evaluating the spread of a fire from a single point in homogenous
conditions. However,
when using Huygen's principle to evaluate Rothermel's model in
heterogeneous spatial conditions, the
fire perimeter must be discretized and evaluated for each point
along the perimeter. This leads to a significant increase in
computational time as the size of the fire increases.
Since the CNN is already analyzing all the input channels as 2-D
image channels, there is no increase in computational cost to
add heterogeneous conditions.
As an
example, a 10km fire perimeter with a 30m grid will require 333
evaluations of the phenomological model, with a total computational
time of 1.0s.
The CNN approach will require a single evaluation, with a total
computational time of 0.004s. 


\begin{figure*}[htbp]
  \centering
	\includegraphics[height=0.16\textheight]{Fig9a1.eps}
	~
	\includegraphics[height=0.16\textheight]{Fig9b1.eps}
	~
	\includegraphics[height=0.16\textheight]{Fig9c1.eps}
	\\
	\includegraphics[height=0.16\textheight]{Fig9a2.eps}
	~
	\includegraphics[height=0.16\textheight]{Fig9b2.eps}
	~
	\includegraphics[height=0.16\textheight]{Fig9c2.eps}
	\\
	\includegraphics[height=0.16\textheight]{Fig9a3.eps}
	~
	\includegraphics[height=0.16\textheight]{Fig9b3.eps}
	~
	\includegraphics[height=0.16\textheight]{Fig9c3.eps}
	\\
	\includegraphics[height=0.16\textheight]{Fig9a4.eps}
	~
	\includegraphics[height=0.16\textheight]{Fig9b4.eps}
	~
	\includegraphics[height=0.16\textheight]{Fig9c4.eps}
	\\
	\includegraphics[height=0.16\textheight]{Fig9a5.eps}
	~
	\includegraphics[height=0.16\textheight]{Fig9b5.eps}
	~
	\includegraphics[height=0.16\textheight]{Fig9c5.eps}
	\\
\hspace{0.022\textwidth}(a)\hspace{0.29\textwidth}(b)\hspace{0.33\textwidth}(c)
\caption{Example CNN prediction of time resolved burn maps (a) Simulation burn maps (b) CNN prediction burn maps (c) Performance metrics at six hour intervals}
\label{fig:timeAnalysis}       % Give a unique label
\end{figure*}

Although the model was trained using a six hour time interval between
the input and output burn maps, 
it is possible to obtain predictions at points
further in the future at six hour intervals by recursively using the previous
prediction as an input to the CNN. Figure~\ref{fig:timeAnalysis} shows
five example cases where this process was used to predict burn maps
up to 24 hours from ignition based on an input burn map six hours after
ignition. The results for each case show the general direction of spread
is captured well, with $F > 0.8$ in all cases except the fourth case where
the input and output fires are small.
The sensitivity of the predictions generally increases
with time, whereas the precision generally decreases with time.
This shows early in the progression of the fire the model
under-predicts the rate of spread, which highlights the difficulty
the CNN can have when dealing with low feature density. Later in
the progression of the fire, the model begins to over-predict the
rate of spread. Since the network was trained on constant rates
of spread for each configuration, the over-prediction is likely
a result of the network over-correcting for the initial
under prediction of rate of spread.































\section{Conclusion}
\label{s:Conclusion}

A novel predictive analytics approach to estimating the spread of
a wildland fire using a convolutional neural network (CNN) was presented.
The robustness of the approach was tested using 1,000 simulations (3,000
burn map pairs) not included when training the network. The predictions
of burn maps from the CNN-based approach agreed with simulation results,
with mean precision, sensitivity, and F-measure of 0.97, 0.92, and 0.93,
respectively across a diverse set of input parameters.
Noise in the input parameters was found to have a minimal affect on
the CNN predictions. The computational cost of the method was found
to be comparable to simulating a phenomological model for homogenous
spatial conditions, and significantly better for heterogeneous spatial
conditions. Although trained on predictions six hours apart, the
CNN-based approach is capable of predicting burn maps further
in the future by recursively using previous predictions as inputs to
the model. The cases where F-measure was observed to be less than 0.80
were found to have input burn maps of less than nine pixels. This
indicates as the fire continues to grow, the predictions will continue
to improve.

This work represents a first step in creating a framework to predict
wildland fire spread based on physics based models and historic data.
Although the data
used to train the CNN in this work was generated using a phenomological
model, the model does not have any information about whether the data
is from a computational fluid dynamics model, phenomological model, or
even experimental measurements of burn maps. Additionally, although
the simulations used homogenous vegetation and
landscapes, the feature learning aspect of the CNN-based method is
well posed to learn heterogenous spatial conditions. The next step
in this process is to incorporate data from simulations with spatially
varying environmental conditions in the training and test sets.








% For one-column wide figures use
%\begin{figure}
% Use the relevant command to insert your figure file.
% For example, with the graphicx package use
%  \includegraphics[width=0.5\textwidth]{example.eps}
% figure caption is below the figure
%\caption{Convolutional Neural Network Architecture}
%\label{fig:1}       % Give a unique label
%\end{figure}

%
%



%\begin{acknowledgements}
%If you'd like to thank anyone, place your comments here
%and remove the percent signs.
%\end{acknowledgements}

% BibTeX users please use one of
%\bibliographystyle{spbasic}      % basic style, author-year citations
%\bibliographystyle{spmpsci}      % mathematics and physical sciences
\bibliographystyle{spphys}       % APS-like style for physics
\bibliography{wildfireReferences}   % name your BibTeX data base

% Non-BibTeX users please use
%\begin{thebibliography}{}
%
% and use \bibitem to create references. Consult the Instructions
% for authors for reference list style.
%
%\bibitem{weber1991}
%Rodney O. Weber, Modelling fire spread through fuel beds, Progress in Energy and Combustion Science, 17, 67-82 (1991)
%\bibitem{rothermel1972}
%Richard C. Rothermel, A mathematical model for predicting fire spread in wildland fuels, U.S. Department of Agriculture, Forest Service, Ogden, UT (1972)
%\bibitem{scott2005}
%Joe H. Scott; Robert E. Burgan, Standard fire behavior fuel models: a comprehensive set for use with Rothermel's surface fire spread model, U.S. Department of Agriculture, Forest Service, Rocky Mountain Research Station, Fort Collins, CO (2005)

% Format for books
%\bibitem{RefB}
%Author, Book title, page numbers. Publisher, place (year)
%\end{thebibliography}

\end{document}

